\section{Legendre history compression}
\label{cp:legendre-section}

\begin{lemma}[Projection identities]
\label{cp:projection}
Let $A>0$ and let $q\ge1$ be an integer. Set
$p_j^A(\xi)=P_j(2\xi/A-1)$, where $P_j$ is the Legendre
polynomial, and let $\Pi_q^A$ project onto degrees below $q$ in
$L^2([0,A])$.  For an absolutely continuous Hilbert-valued history $u$,
\begin{align}
 \|(I-\Pi_q^A)u\|_{L^2}^2
 &\le\frac1{q(q+1)}\int_0^A\xi(A-\xi)\|u'(\xi)\|^2\,d\xi,
 \label{cp:eq-projection-tail}\\
 \|u(A)-(\Pi_q^Au)(A)\|
 &\le\sqrt{A/(3q)}\,\|u'\|_{L^2},&
 \|(\Pi_q^Au)(A)\|&\le64\sqrt q\,\|u\|_{L^\infty}.
 \label{cp:eq-projection-endpoint}
\end{align}
For a fixed absolutely continuous history $g$, the growing-interval error obeys
\begin{equation}
 \frac d{dA}\int_0^A\|g-\Pi_q^Ag\|^2\,d\xi
       =\|g(A)-(\Pi_q^Ag)(A)\|^2.
 \label{cp:eq-projection-growth}
\end{equation}
\end{lemma}
\begin{proof}
Rodrigues' formula and integration by parts give
$\int_0^Ap_i^Ap_j^A=A\mathbf1_{i=j}/(2j+1)$ and
$-[\xi(A-\xi)(p_j^A)']'=j(j+1)p_j^A$.
The derivative polynomials are orthogonal for the weight
$\xi(A-\xi)$.  Bessel's inequality followed by the ordinary expansion
therefore proves \eqref{cp:eq-projection-tail}; polynomial approximation
and increasing finite-dimensional projections give the stated
absolutely continuous, Hilbert-valued version.
Integration against each lower polynomial also gives
$P_j'=\sum_{i<j,\ j-i\ \mathrm{odd}}(2i+1)P_i$.  Thus the endpoint kernel is
\[
 K_q^A(A,\xi)=\frac1A\sum_{j<q}(2j+1)p_j^A(\xi)
            =\tfrac12[(p_q^A)'(\xi)+(p_{q-1}^A)'(\xi)].
\]
Its integral is one, and integration by parts gives
$u(A)-(\Pi_q^Au)(A)=\tfrac12\int_0^A(p_q^A+p_{q-1}^A)u'$.
The squared $L^2$ norm of this multiplier is
$A[(2q+1)^{-1}+(2q-1)^{-1}]/4\le A/(3q)$.

For completeness the needed kernel $L^1$ estimate has a short
one-dimensional proof.  Put
$v(\theta)=\sqrt{\sin\theta}P_k(\cos\theta)$ and
$a_k=(k+1/2)^2+(4\sin^2\theta)^{-1}$.  The Legendre equation implies
$v''+a_kv=0$, whence
$\partial_\theta(v^2+v'^2/a_k)=-a_k'v'^2/a_k^2\ge0$ on
$(0,\pi/2)$.  At the midpoint use
$P_{2j}(0)=(-1)^j\binom{2j}j/4^j$,
$P_{2j+1}'(0)=(2j+1)P_{2j}(0)$, and
$(\binom{2j}j/4^j)^2\le1/(j+1)$; the last inequality follows by
induction.  The midpoint energy is at most $4/k$.  Differentiating
$P_k=v/\sqrt{\sin\theta}$ now yields
\[
 |\partial_\theta P_k(\cos\theta)|
 \le3\sqrt{k}(\sin\theta)^{-1/2}
       +2k^{-1/2}(\sin\theta)^{-3/2}.
\]
On $[1/k,\pi/2]$ its integral is at most $3\pi\sqrt{k}+8$,
using $\sin\theta\ge2\theta/\pi$.  On $[0,1/k]$ use
$|P_k'|\le k(k+1)/2$ to obtain at most $1/2$.
This derivative bound follows from the derivative expansion and
$|P_j|\le1$, the latter following from the integral formula
$P_j(\cos\theta)=\pi^{-1}\int_0^\pi
(\cos\theta+i\sin\theta\cos\psi)^j\,d\psi$ (expand its generating
function).  Reflection gives total variation at most $36\sqrt k$.
The endpoint kernel consequently has $L^1$ norm at most
$36\sqrt q\le64\sqrt q$; $q=1$ is the averaging kernel.
Finally differentiation of the squared projection error proves
\eqref{cp:eq-projection-growth}: the interior term vanishes because
$\partial_A\Pi_q^Ag$ is still a polynomial of degree below $q$.
\end{proof}

\begin{proposition}[Legendre approximation at the prescribed size]
\label{cp:legendre}
Eventually at each fixed confidence, there is an initialization-only
Legendre model with
\begin{equation}
 q=\left\lceil C\left(\frac m\gamma\right)^3n^{1/4}
             [\log(en)]^2e^{\sqrt{\log(en)}/2}\right\rceil
 \label{cp:eq-legendre-order}
\end{equation}
such that
\begin{equation}
 \sup_{t\in[0,\infty],\ \|x\|=\sqrt d}
       |f_{\rm Leg}(t,x)-f_n(t,x)|
 \le\frac{CY}{\sqrt n[\log(en)]^3}.
 \label{cp:eq-legendre-accuracy}
\end{equation}
Its moving storage is at most
\[
 Cm(m/\gamma)^3n^{5/4}[\log(en)]^2e^{\sqrt{\log(en)}/2};
\]
the fixed storage is the additional $(L-1)n^2$.
\end{proposition}

\begin{proof}
\noindent\textbf{Local notation and construction.}\quad
All auxiliary quantities below belong to this proof. Put
\[
 \lambda=\gamma/m,\qquad
 H=\max_{0\le j\le L}\sqrt{Q^{(j)}_{11}},\qquad
 s=\max\left\{1,\max_j\sup_{|\operatorname{Im}z|\le a/2}
                     |\phi_j'(z)|\right\},\qquad
 b=\max_j|\phi_j(0)|,
\]
and define $D_j=s(9s)^{L-j}$. These definitions are made here;
no coefficient from the proof of \Cref{cp:fit} is imported.
In particular $H\le(2\beta)^L$, by the affine population-RMS recursion,
and $s,b\le\beta$.

The Legendre model constructed here shares the realized dense initialization and physical
time.  Here the whole-sphere version of \eqref{cp:norm} is abbreviated by
\[
 \|g\|_*=\sup_{t\in[0,\infty],\,\|x\|=\sqrt d}|g(t,x)|,
\]
including the fitted limit at $t=\infty$.
Write $\widehat f=f_{\rm Leg}=n^{-1}\widehat w^\top\widehat h^{(L)}$,
$\widehat r_a=\widehat f(x_a)-y_a$ and
$\widehat\rho=\|\widehat r\|_2/\sqrt m$. We first construct the model for
an arbitrary integer $q\ge1$ and then take the order in the statement.
At each hidden interface
$\ell=2,\ldots,L$, sample $a$, and order $j=0,\ldots,q-1$, retain
$\bar h_{a,j}^{(\ell-1)},\bar\delta_{a,j}^{(\ell)}\in\mathbb R^n$.
Together with $\widehat W^{(1)}$, $\widehat w$, and a scalar $\tau$,
these are the moving state.  Reconstruct
\begin{equation}
 \widehat W^{(\ell)}=W_0^{(\ell)}-
 \frac2{mn\tau}\sum_{a=1}^m\sum_{j=0}^{q-1}(2j+1)
       \bar\delta_{a,j}^{(\ell)}\bar h_{a,j}^{(\ell-1)\top}.
 \label{cp:eq-legendre-reconstruction}
\end{equation}
Use these matrices in the ordinary forward pass and in
$\widehat k_a^{(L)}=\widehat w$,
$\widehat\delta_a^{(\ell)}=
\phi_\ell'(\widehat z_a^{(\ell)})\odot\widehat k_a^{(\ell)}$,
$\widehat k_a^{(\ell)}=
\widehat W^{(\ell+1)\top}\widehat\delta_a^{(\ell+1)}$.
The autonomous physical-time equations are
\begin{align}
 \dot{\widehat W}^{(1)}&=-\frac2m\sum_a
     \widehat r_a\widehat\delta_a^{(1)}v_a^\top,&
 \dot{\widehat w}&=-\frac2m\sum_a\widehat r_a\widehat h_a^{(L)},&
 \dot\tau&=\widehat\rho,\nonumber\\
 \dot{\bar h}_{a,j}^{(\ell-1)}
 &=\widehat\rho\widehat h_a^{(\ell-1)}-
 \frac{\widehat\rho}{\tau}
 \left(j\bar h_{a,j}^{(\ell-1)}+
       \sum_{i<j}(2i+1)\bar h_{a,i}^{(\ell-1)}\right),\nonumber\\
 \dot{\bar\delta}_{a,j}^{(\ell)}
 &=\widehat r_a\widehat\delta_a^{(\ell)}-
 \frac{\widehat\rho}{\tau}
 \left(j\bar\delta_{a,j}^{(\ell)}+
       \sum_{i<j}(2i+1)\bar\delta_{a,i}^{(\ell)}\right).
 \label{cp:eq-legendre-flow}
\end{align}
Initially $\tau=1$, $\widehat W^{(1)}=W_0^{(1)}$, $\widehat w=0$,
$\bar h_{a,0}^{(\ell-1)}=h_{0,a}^{(\ell-1)}$, and all other moments
vanish.  No equation divides by a residual.  At $\widehat\rho=0$ all
velocities vanish, and the vector field is locally Lipschitz for
$\tau>0$.  The exact inventory is
\begin{equation}
 S_{\rm moving}=n(d+1)+1+2(L-1)mnq,\qquad
 S_{\rm fixed}=(L-1)n^2.
 \label{cp:eq-legendre-count}
\end{equation}
The latter entries are the initialized mixers.  Data storage is
additional; reconstructed matrices need not be retained independently.

\noindent\textbf{Moment identities and the physical defect.}\quad
For use in this proof, let $p_j^A(\xi)=P_j(2\xi/A-1)$ and let
$\Pi_q^A$ be the degree-below-$q$ orthogonal projection in $L^2([0,A])$,
as in \Cref{cp:projection}.
The histories in the proof are specified by
$h_a(\tau(t))=\widehat h_a(t)$ and
$b_a(\tau(t))=(\widehat r_a/\widehat\rho)\widehat\delta_a(t)$.
On the prefix $[0,1]$ they equal the initialized forward feature and
zero, respectively.  The backward history is continuous at the join
because the initial readout is zero.  The moments in
\eqref{cp:eq-legendre-flow} are exactly $\int_0^\tau p_j^\tau h_a$
and $\int_0^\tau p_j^\tau b_a$; differentiating these integrals proves
the equations, including their factors $j$ and $2i+1$.
Differentiating their bilinear projection pairing gives the exact defect
\begin{equation}
 \dot{\widehat W}^{(\ell)}
 =-\frac2{mn}\sum_a\widehat r_a\widehat\delta_a^{(\ell)}
                       \widehat h_a^{(\ell-1)\top}+\mathcal E_\ell,
 \qquad
 \mathcal E_\ell=\frac{2\widehat\rho}{mn}
       \sum_a e_{b,a}^{(\ell)}e_{h,a}^{(\ell-1)\top},
 \label{cp:eq-legendre-defect}
\end{equation}
where $e_g=g(\tau)-(\Pi_q^\tau g)(\tau)$.
Indeed the derivative of the pairing of the two projected histories
equals the full endpoint product minus the product of their endpoint
errors, by orthogonality.  These identities hold until a residual zero
and extend there by stationary continuation.  Histories are proof
objects, not extra state.

\begin{claim}[Order-independent physical bounds]
\label{cp:legendre-fit}
On \eqref{cp:eq-initialization-event}, under
$Y\le\lambda\beta^{-30L}$, every finite order is globally defined,
has convergent physical parameters and moments, and fits the labels.
Its operator and feature bounds are those of
\eqref{cp:eq-real-tube}, and
\begin{equation}
 \widehat\rho(t)\le Ye^{-\lambda t/4},\qquad
 \tau(t)\le1+4Y/\lambda,\qquad
 \|\widehat w(t)\|/\sqrt n\le4Y/\sqrt\lambda.
 \label{cp:eq-legendre-physical}
\end{equation}
\end{claim}
\emph{Verification of the local claim.}
Let $\sigma=8Y/\lambda$, $R=\sigma\sqrt\lambda$, and
$\kappa=\lambda/4$. Define, within the present construction proof,
\[
 C_1=4s^2D_1^2,\quad
 C_\ell=3s^2(64H^4D_\ell^2+82C_{\ell-1}),\quad
 E_*=130\sum_{\ell=2}^LD_\ell\sqrt{C_{\ell-1}}.
\]
It suffices that
\begin{equation}
 \sigma\le\min\left\{1,\frac1{64H^2D_1},
 \frac1{\sqrt{8\sqrt{C_L}}},
 \left(\frac1{8H^2D_1E_*}\right)^{2/7}\right\}.
 \label{cp:eq-legendre-fitting-smallness}
\end{equation}
The present label cap implies this: $D_\ell\le\beta^{2L}$,
$\sqrt{C_\ell}\le\beta^{8L}$, and $E_*\le\beta^{12L}$.
For the middle estimate unroll
$C_\ell=192s^2H^4D_\ell^2+246s^2C_{\ell-1}$, bound
$192s^2,246s^2\le\beta^5$, $H^4D_\ell^2\le\beta^{10L}$,
and sum at most $L$ terms.  For the last use
$130L\le(\beta^L)^2$.  Substitution gives a lower bound
$\beta^{-8L}$ for the minimum in
\eqref{cp:eq-legendre-fitting-smallness}, whereas
$\sigma\le8\beta^{-30L}$.

Stop at clock activity $\sigma$, readout RMS $R$, operator norm nine,
or feature RMS $2H$.  Before the stop,
$\|\widehat\delta_a^{(\ell)}\|/\sqrt n\le D_\ell R$.
Projection contraction, sample Cauchy--Schwarz, and the zero backward
prefix give
\[
 \|\widehat W^{(\ell)}-W_0^{(\ell)}\|_F
 \le4HD_\ell R\sqrt{(1+\sigma)\sigma},\qquad
 \|\widehat W^{(1)}-W_0^{(1)}\|_F/\sqrt n\le2D_1R\sigma.
\]
These are strictly below one by
\eqref{cp:eq-legendre-fitting-smallness} and $\sqrt\lambda\le H$.
Let $Z_\ell$ be the supremum, over queries and stopped horizons, of
$n^{-1}\int_1^{\tau(t)}\|\partial_\xi h^{(\ell)}(\xi,v)\|^2d\xi$.
The polynomial endpoint norm on $L^2([0,A])$ is $q/\sqrt A$,
since the sum of squared normalized basis values at $A$ is
$\sum_{j<q}(2j+1)/A=q^2/A$. Thus the aggregate backward endpoint
error is at most $(q+1)D_\ell R$.  By
\eqref{cp:eq-projection-growth}, \eqref{cp:eq-projection-tail}, and
\eqref{cp:eq-legendre-defect},
\[
 \int_0^t\widehat\rho
       \|\mathcal E_\ell/\widehat\rho\|_F^2du
 \le2(1+\sigma)^2D_\ell^2R^2Z_{\ell-1}.
\]
Here the forward prefix error is zero,
$\xi(A-\xi)\le A^2/4$, and $(q+1)/q\le2$.
The first-layer chain rule and the three later derivative terms
(gradient update, defect, and propagated feature derivative) now give
\[
 Z_1\le4\sigma s^2D_1^2R^2,\qquad
 Z_\ell\le3s^2\{64\sigma H^4D_\ell^2R^2+
                 [81+32H^2D_\ell^2R^2]Z_{\ell-1}\}
          \le C_\ell\sigma R^2.
\]
The last step uses $32H^2D_\ell^2R^2<1$.
Consequently every feature moves in RMS by at most
$\sqrt{\sigma Z_\ell}\le\sqrt{C_L}\sigma^2\sqrt\lambda
\le\sqrt\lambda/8$.  This excludes the feature boundary and preserves
the top Gram gap $\lambda/4$ by singular-value perturbation.

For the pointwise defect use the sharper bounds
\eqref{cp:eq-projection-endpoint}.  The two endpoint errors are at most
$65\sqrt qD_\ell R$ and
$\sqrt{(1+\sigma)/(3q)}\sqrt{Z_{\ell-1}}$, respectively.  Hence
\begin{equation}
 \sum_{\ell=2}^L\|\mathcal E_\ell\|_F
          \le E_*\sqrt\sigma R^2\widehat\rho.
 \label{cp:eq-legendre-pointwise-defect}
\end{equation}
The prediction Jacobian applied to this hidden defect has sample RMS
at most $2HD_1R\sum\|\mathcal E_\ell\|_F
\le\lambda\widehat\rho/4$, by the last entry in
\eqref{cp:eq-legendre-fitting-smallness}. Let $J$ be the Jacobian of
the training prediction vector in mobility coordinates
$(W^{(1)}/\sqrt n,W^{(2)},\ldots,W^{(L)},w/\sqrt n)$, evaluated at
the reconstructed state, and put $\Gamma=JJ^\top/m$.
The mean full tangent Gram then gives
$\dot{\widehat r}=-2\Gamma\widehat r+J\mathcal E$ with
$\Gamma\succeq\lambda I/4$.  Thus
$\dot{\widehat\rho}\le-\kappa\widehat\rho$ and
$\int_t^T\widehat\rho\le\widehat\rho(t)/\kappa$ on every stopped
interval.  The clock activity is at most $\sigma/2$.

Let $\mathcal F$ be the dense gradient vector field at the reconstructed
parameters, in the norm of Lemma~\ref{cp:fit}.  Its squared norm is
$4\widehat r^\top\Gamma\widehat r/m\ge\lambda\widehat\rho^2$,
and the defect estimate implies
$-\partial_t\widehat\rho^2\ge\|\mathcal F\|_{\rm par}^2/2$.
Multiplying by $e^{\kappa t}$, integrating by parts, and using
$\widehat\rho\le Ye^{-\kappa t}$ gives
$\int e^{\kappa t}\|\mathcal F\|_{\rm par}^2dt\le4Y^2$.
Weighted Cauchy--Schwarz then gives
$\int\|\mathcal F\|_{\rm par}\le2Y/\sqrt\kappa=R/2$.
The readout has no defect, so its boundary is excluded as well.
Every stopped inequality improves strictly.  The integral formulas
bound all moments at finite $n,q$, proving global continuation.
Uniqueness excludes a finite first residual zero when $Y>0$, since
zero residual is an equilibrium of the entire stored-state ODE.
The integrable gradient and defect lengths yield physical limits;
the moment equations also have integrable velocities.  This proves
\eqref{cp:eq-legendre-physical} and fitting at the limit.
This finishes the local fitting claim. The coefficients $C_j$
just defined remain private to this construction proof.

\noindent\textbf{Comparison with dense training.}\quad
Intersect the initialization event and the all-time training-carrier
bound \eqref{cp:carrier} of \Cref{cp:source}. Set
$X=\beta^L$, $z=Y/\lambda$, $\chi=1+\lambda^{-1}$,
$u=\sqrt{\log(en)}$, and $M=32X^{21}zu$.
Write $\rho_D=\|r\|_2/\sqrt m$ for the dense residual RMS.
For the backward comparison at any parameter state, define its pre-gated
training responses by $k_a^{(L)}=w$ and
$k_a^{(j)}=W^{(j+1)\top}\delta_a^{(j+1)}$ for $j<L$.
Only actual dense training carriers, not closure carriers or passive
query carriers, are bounded by $M$.  Throughout this proof
$z\le X^{-30}$, $X\ge100$. Use the Euclidean mobility coordinates
$(W^{(1)}/\sqrt n,W^{(2)},\ldots,W^{(L)},w/\sqrt n)$,
whose Euclidean norm is the parameter norm of \Cref{cp:fit}.
Let $e$ be the Euclidean discrepancy and $D$ the
sum of its $L+1$ block norms, so $D\le\sqrt{L+1}e$.

We first verify the signed comparison with no reciprocal-gap factor
inserted into its Jacobian coefficient.  Put $T_0=9s$,
$R=8z\sqrt\lambda$, $B_\delta=sT_0^{L-1}R$, and define
\begin{gather*}
 Q_1=b+9s,\quad Q_\ell=b+9sQ_{\ell-1},\quad
 Q=\max(1,Q_1,\ldots,Q_L),\\
 F_{\rm seg}=QT_0^{L-1},\quad P=1+Q(L-1),\\
 d_0=sT_0^{L-1}(1+B_\delta),\quad
 d_1=Lt_2F_{\rm seg}T_0^{L-1},\\
 K_0=\sqrt{L+1}\{Pd_0+[1+(L-1)B_\delta]sF_{\rm seg}\},
 \qquad K_1=\sqrt{L+1}Pd_1,
\end{gather*}
where $t_2=\max\{1,\max_\ell\sup_{|\operatorname{Im}z|\le a/2}
|\phi_\ell''(z)|\}\le\beta$.
On the straight parameter segment from dense to closure, matrix and
readout norms keep their bounds by convexity; linear growth gives
feature RMS at most $Q$.  Forward subtraction therefore bounds the
preactivation difference at segment fraction $\vartheta\in[0,1]$ by $F_{\rm seg}\vartheta D$.
In backward subtraction use
\[
 \delta_\vartheta-\delta_D
 =\phi'(z_\vartheta)\odot(k_\vartheta-k_D)
       +[\phi'(z_\vartheta)-\phi'(z_D)]\odot k_D.
\]
Propagation costs $T_0$, a changed mixer acts on a response of RMS at
most $B_\delta$, and the changed gate uses only the dense carrier $M$.
Induction gives response difference at most $(d_0+d_1M)\vartheta D$.
Writing $f_a(\theta)=f_\theta(x_a)$, the gradient blocks of a training output are
$\delta_a^{(1)}v_a^\top/\sqrt n$,
$\delta_a^{(\ell)}h_a^{(\ell-1)\top}/n$, and
$h_a^{(L)}/\sqrt n$.
Subtract each hidden product as
$(\delta_\vartheta-\delta_D)h_\vartheta^\top+\delta_D(h_\vartheta-h_D)^\top$.
Their block bounds give
$\|\nabla f_a(\theta_\vartheta)-\nabla f_a(\theta_D)\|
\le(K_0+K_1M)\vartheta e$.
Segment integration proves both hypotheses of Lemma~\ref{cp:signed}
with $K=K_0+K_1M$.  The residual integrals are at most $2z$ and $4z$,
so for every finite horizon
\begin{equation}
 \sup_{t\le T}D(t)\le\mathcal A_n\mathcal D_T,
 \quad\mathcal A_n=\sqrt{L+1}e^{14(K_0+K_1M)z},
 \quad\mathcal D_T=\int_0^T\sum_{\ell=2}^L\|\mathcal E_\ell\|_Fdt.
 \label{cp:eq-legendre-signed}
\end{equation}
In particular the carrier contribution is additive in $K$, not
multiplied once per layer.

Here are coefficient bounds needed to control the forcing.  Their
dependence on the gap and labels is displayed before any simplification.
Put $a_0=4z$, $\kappa=\lambda/4$, $F_z=2HT_0^{L-1}$,
$P_0=1+2H(L-1)$, retain the local coefficients $C_\ell$ defined above, and define
\begin{align*}
 B_d&=sT_0^{L-1}(1+B_\delta)+Lt_2F_zT_0^{L-1}M,\\
 V_z&=2F_zB_\delta P_0,\\
 V_\delta&=T_0^{L-1}
   [4sH+4sH(L-1)B_\delta^2+Lt_2MV_z],\\
 T_1&=2D_1R,\qquad
 T_\ell=4HD_\ell R+130D_\ell\sqrt{8zC_{\ell-1}}R^2,\\
 V_1&=sT_1,\qquad V_\ell=s(2HT_\ell+9V_{\ell-1}),
       \qquad V_h=\max_\ell V_\ell,\\
 C_c&=4\left(4H^2+D_1^2R^2+
                    4H^2R^2\sum_{\ell=2}^LD_\ell^2\right)+\lambda/2,\\
 F_h&=V_ha_0\sqrt{(1+a_0)/2},\\
 b_0&=\sqrt{1+a_0}(Y/\kappa)V_\delta+2B_\delta\sqrt{a_0},
       \qquad b_1=2\sqrt{1+a_0}\,C_cB_\delta/\kappa,\\
 C_f&=2H+RsF_z.
\end{align*}
The meanings follow directly from the recursions: $B_dD$ bounds the
response discrepancy between the two actual trajectories; $V_z\rho_D$
and $V_\delta\rho_D$ bound dense preactivation and response speeds;
$V_h\widehat\rho$ bounds closure feature speeds, using
\eqref{cp:eq-legendre-pointwise-defect}.  Finally
$\|\partial_t(\widehat r/\widehat\rho)\|_m\le C_c$ follows by
differentiating the normalized residual and bounding its tangent Gram
by its trace and its defect by $\lambda\widehat\rho/4$.

We spell out sufficient power bounds to fix the exponential rate.
The preceding finite recursions imply
\begin{align}
 &K_0\le X^9,\quad K_1\le X^{13},\quad
 14(K_0+K_1M)z\le X^{10}z+X^{36}z^2u,\nonumber\\
 &V_h\le X^{20}R,\quad C_c\le X^{11},\quad
 B_d\le X^{33}(1+zu),\quad V_\delta\le X^{40}(1+zu),\nonumber\\
 &F_h/Y\le X^{22}z\sqrt\chi,\quad
 b_0\le X^{42}z(1+zu),\quad b_1\le X^{16}z\sqrt\chi,
 \quad C_f\le X^7.
 \label{cp:eq-legendre-ledger}
\end{align}
To check the first line, $Q\le X^2$, $F_{\rm seg}\le X^4$,
$d_0\le2X^3$, $d_1\le X^8$, $P\le X^4$ give the first two
bounds; insert $M\le32X^{21}zu$ and use $448\le X^2$.
For the second line, $R\le X^{-27}$, $\sqrt{C_\ell}\le X^8$,
and $D_\ell\le X^2$ give $T_\ell\le X^{12}R$.
The $V_\ell$ recursion has forcing at most $2X^{15}R$, propagation
at most $T_0^L\le X^2$, and at most $L\le X$ terms.
The trace defining $C_c$ gives its stated bound.
Also $F_z\le X^5$, $P_0\le X^4$, and $V_z\le2X^{12}R$.
The non-carrier and carrier terms of $V_\delta$ are at most
$4LX^5$ and $64LX^{36}Rzu$, respectively; those of $B_d$ are
at most $2X^3$ and $32LX^{29}zu$.
For the final line retain the normalization:
$F_h/Y\le32X^{20}z/\sqrt\lambda$,
$b_1\le96X^{14}z/\sqrt\lambda$, and
$b_0\le6X^{40}z(1+zu)+32X^5z^{3/2}$.
These inequalities imply \eqref{cp:eq-legendre-ledger} because
$X\ge100$, $z\le1$, and $\lambda^{-1/2}\le\sqrt\chi$.
In particular
\begin{equation}
 \mathcal A_n\le eX e^u,
 \qquad \log\mathcal A_n=O_{\rm problem}(\sqrt{\log(en)}).
 \label{cp:eq-legendre-exponent}
\end{equation}
The universal coefficient of $u$ follows from $X^{36}z^2\le1$;
no extra power of $m/\gamma$ occurs in that exponent.

By \eqref{cp:eq-projection-tail}, the forward history tail on
$[0,A]$, $A\le1+a_0$, has aggregate RMS $L^2$ norm at most $F_h/q$:
use its clock derivative bound $V_h$ and
$\int_1^A\xi(A-\xi)d\xi\le A(A-1)^2/2$.
For the backward history compare it with
$\widetilde b_a(\tau(t))=
(\widehat r_a/\widehat\rho)\delta_{D,a}(t)$, with zero prefix.
Its physical derivative has aggregate RMS at most
$C_cB_\delta+V_\delta\rho_D$.
Freeze it after $t_q=2\log q/\kappa$ (or the current horizon if
earlier).  Since $A-\tau(t)\le\widehat\rho(t)/\kappa$, the weighted
derivative energy cancels the inverse clock speed and is at most
\[
 \frac{1+a_0}{\kappa}
 \left[2t_qC_c^2B_\delta^2+
       2V_\delta^2\int_0^{t_q}\rho_D(t)^2dt\right].
\]
The frozen tail is therefore at most
$(b_1\sqrt{\log q}+\sqrt{1+a_0}(Y/\kappa)V_\delta)/q$.
The freezing error is at most
$2B_\delta\sqrt{a_0}e^{-\kappa t_q/2}
=2B_\delta\sqrt{a_0}/q$.
The difference between closure and dense response histories contributes
at most $B_d\sqrt{a_0}\sup_{t\le T}D(t)$; the normalized residual has
sample norm one.  All aggregate history norms here use the factor
$(mn)^{-1/2}$.

Use Cauchy--Schwarz in samples and then in $dA=\widehat\rho\,dt$
in \eqref{cp:eq-legendre-defect}.  The initial projection errors vanish,
so \eqref{cp:eq-projection-growth} gives
\begin{equation}
 \mathcal D_T\le2(L-1)F_h\left[
 \frac{b_0+b_1\sqrt{\log q}}{q^2}
 +\frac{B_d\sqrt{a_0}\sup_{t\le T}D(t)}q\right].
 \label{cp:eq-legendre-absorption}
\end{equation}
For
$q\ge q_{\rm abs}:=4(L-1)F_hB_d\sqrt{a_0}\mathcal A_n$,
\eqref{cp:eq-legendre-signed} absorbs the last term with factor
at most one half.  Whole-sphere forward/readout subtraction costs
$C_fD$.  The coefficients do not depend on $T$, so increasing horizons
and the already constructed physical limits yield
\[
 \|f_{\rm Leg}-f_n\|_*
 \le\frac{4(L-1)C_f\mathcal A_nF_h}{q^2}
                         (b_0+b_1\sqrt{\log q}).
\]
By the ledger its numerator divided by $Y$ is at most
$eX^{75}z^2\chi(1+zu)e^u\sqrt{\log(eq)}$.
The cap and $\lambda\le H^2\le X^4$ imply $Y\le X^{-26}\le1$,
so $z\le\lambda^{-1}\le\chi$.  Consequently
\begin{equation}
 \|f_{\rm Leg}-f_n\|_*
 \le C Y\chi^6 e^{\sqrt{\log(en)}}
       \frac{\sqrt{\log(en)\log(eq)}}{q^2}\qquad(q\ge q_{\rm abs}).
 \label{cp:eq-legendre-forward}
\end{equation}
This enlargement preserves the displayed sample/gap power without
spending the label cap to reduce it.  All coefficients except
$\mathcal A_n$ grow at most polynomially in $u$, hence
$q_{\rm abs}=n^{o(1)}$ for the fixed problem.  The prescribed order
\eqref{cp:eq-legendre-order} eventually exceeds it.  Since
$m/\gamma\ge1$, $\chi\le2m/\gamma$, and
$\log(eq)\le C\log(en)$ eventually, substitution proves
\eqref{cp:eq-legendre-accuracy}.  The count follows from
\eqref{cp:eq-legendre-count}; $nd$ and rounding are absorbed only after
taking sufficiently large width for the fixed problem.  One event
supports every order used here, so no order-dependent failure union is
needed.  All initialization data are the specified dense initialization
and initialized training features.
\end{proof}

\section{An actual dense-trajectory fluctuation}
\label{cp:dense-lower-section}

\begin{lemma}[Derivative to trajectory]
\label{cp:derivative-transfer}
Let $r>0$ and $M\ge0$. If the scalar function $g$ is holomorphic on a
neighborhood of the filled parameter-two
Bernstein ellipse for $[0,r]$ and bounded there by $M$, then for every
integer $N\ge1$,
\begin{equation}
 \sup_{0\le t\le r}|g(t)|
       \ge\frac r{2N^2}|g'(0)|-24M2^{-N}.
 \label{cp:eq-derivative-transfer}
\end{equation}
\end{lemma}
\begin{proof}
Set $t=r(x+1)/2$.  The Joukowski map
$x=(z+z^{-1})/2$ and Cauchy's coefficient formula give Chebyshev
coefficients $|a_k|\le2M2^{-k}$.  The degree-$N$ truncation has
uniform tail at most $2M2^{-N}$; its physical derivative tail at zero
is at most
$4M r^{-1}2^{-N}(N^2+4N+6)$, by summing
$\sum_{k>N}k^22^{-k}$.
For a degree-$N$ polynomial, including complex coefficients,
$|P'(-1)|\le N^2\|P\|_{[-1,1]}$.
Indeed interpolate at $x_j=\cos(j\pi/N)$: the endpoint derivative
weights have alternating signs, and their absolute sum equals
$|T_N'(-1)|=N^2$ by evaluating on $T_N(x_j)=(-1)^j$.
The sign assertion follows directly from the products defining the
Lagrange basis; the endpoint weight is
$\sum_{j<N}(-1-x_j)^{-1}<0$.
Apply this inequality to the truncation, include the physical factor
$2/r$, and add the two tails.  Rearranging gives error coefficient
$4+8/N+12/N^2\le24$, as required.
\end{proof}

\begin{proposition}[Positive-time dense variability]
\label{cp:dense-lower}
Under the assumptions of \Cref{cp:setup}, for every fixed $0<\delta<1$,
there is a constant
$c_{\phi,L,\delta}>0$, depending only on the activations, depth and
this confidence, such that all sufficiently large individual widths satisfy
\begin{equation}
 \Pr\left\{\|f_n-\widetilde f_n\|_{\mathcal X}
       \ge c_{\phi,L,\delta}\frac{Y\sqrt\gamma}
              {\sqrt n[\log(en)]^{5/2}}\right\}\ge1-\delta.
 \label{cp:eq-dense-lower}
\end{equation}
This holds for either the whole sphere or any declared panel containing
the training inputs.  The witness is at a positive training time;
no endpoint lower bound is asserted.
\end{proposition}

\begin{proof}
All moments and fluctuation variables introduced here are private to this
proof. Let $G\sim N(0,1)$, define
\[
 \mu_{2,j}=Q^{(j)}_{11}\quad(j=L-1,L),\qquad
 \mu_4=\mathbb E\phi_L(\sqrt{\mu_{2,L-1}}G)^4,
\]
and write $\Phi$ for the standard normal distribution function.
Let $Z\sim N(0,Q^{(L-1)})$, $U_a=\phi_L(Z_a)$ and
$Q=\mathbb E UU^\top=Q^{(L)}$. We first establish the innovation and
conditional fluctuation needed below, then convert an initial derivative
into a positive-time trajectory separation.

\begin{claim}[A final-layer innovation]
\label{cp:innovation}
For $m\ge2$ and $Y>0$, some deterministic training index $a_*$ satisfies
\begin{equation}
 \operatorname{Var}\{U_{a_*}(U^\top y)\}
       \ge\frac{\gamma^3\mu_{2,L}}{16\mu_4}Y^2>0.
 \label{cp:eq-innovation}
\end{equation}
\end{claim}
\emph{Verification.}
All fourth moments are finite by linear growth.  The positive gap gives
$\mu_{2,L}>0$ and $\mu_4>0$.  Put $c=Qy$, $A=\|c\|$,
$T=\operatorname{tr}Q$, $M_4=\mathbb E\|U\|^4$, and
$D=T-c^\top Qc/A^2>0$.  The last inequality holds because $m\ge2$
and $Q\succeq\gamma I$.  For $V=\mathbb E\|U(U^\top y)-c\|^2$,
the squared-area identity gives
\[
 A^2\|U\|^2-(c^\top U)^2
 \le\|U(U^\top y)-c\|^2\|U\|^2.
\]
Splitting at $\|U\|=R_0$ yields
$A^2D\le R_0^2V+A^2M_4/R_0^2$.
Choose $R_0^2=2M_4/D$ to obtain
$V\ge[y^\top Q^2(TI-Q)y]^2/(4M_4\|Qy\|^2)$.
Here $M_4\le m^2\mu_4$ and $T=m\mu_{2,L}$.
Every eigenvalue $x$ of $Q$ lies in
$[\gamma,T-(m-1)\gamma]$.  Thus, for
$N=y^\top Q^2(TI-Q)y$,
\[
 N\ge(m-1)\gamma\|Qy\|^2,\qquad
 N\ge\gamma^2(T-\gamma)\|y\|^2.
\]
The second uses
$x(T-x)-\gamma(T-\gamma)=(x-\gamma)(T-\gamma-x)\ge0$
and then $x\ge\gamma$.  Multiplying these lower bounds gives
\[
 V\ge\frac{(m-1)\gamma^3(m\mu_{2,L}-\gamma)}{4m^2\mu_4}\|y\|^2
 \ge\frac{\gamma^3\mu_{2,L}}{16\mu_4}\|y\|^2,
\]
because $\gamma\le\mu_{2,L}$ and $(m-1)^2/m^2\ge1/4$.
At least one of the $m$ coordinate variances is at least $V/m$,
which proves \eqref{cp:eq-innovation}.

\begin{claim}[Conditional scalar fluctuation]
\label{cp:scalar-fluctuation}
For the index in the preceding local claim, put
$g_n(t)=f_n(t,x_{a_*})-\widetilde f_n(t,x_{a_*})$, where the two
dense initializations are independent.  There is a fixed $s_*>0$ with
\begin{equation}
 s_*^2\ge\frac{\gamma^3\mu_{2,L}Y^2}{2m^2\mu_4},\qquad
 \liminf_{n\to\infty}\Pr\{\sqrt n|g_n'(0)|\ge u s_*\}
       \ge2[1-\Phi(u)]\quad(u>0),
 \label{cp:eq-scalar-fluctuation}
\end{equation}
where $\Phi$ is the standard normal distribution function.
\end{claim}
\emph{Verification.}
Let $\mathcal F_n$ contain both initializations through layer $L-1$.
Their training empirical Grams converge to $Q^{(L-1)}$ by
\eqref{cp:eq-initial-gram-limit} of \Cref{cp:fit}, including when this
matrix is singular.
Conditionally on $\mathcal F_n$, the last-layer preactivation rows of
each run are independent Gaussian vectors. Applying $\phi_L$ gives
independent feature vectors, generally non-Gaussian; the two groups are
also independent.
For either run write $U_{i,a}=h_{0,i}^{(L)}(v_a)$ for its initialized
top-layer feature entries. The zero readout gives exactly
\[
 \dot f_n(0,x_a)=\frac2{mn}\sum_{i=1}^n U_{i,a}(U_i^\top y).
\]
Write $v_* =\operatorname{Var}[U_{a_*}(U^\top y)]>0$ and
$s_*^2=8v_*/m^2$.  Conditional variances of the row products converge
to $v_*$: couple their Gaussian vectors with the continuous covariance
square root and dominate the fourth-degree products by Gaussian
polynomials.  Their centered absolute third moments are uniformly
bounded on bounded covariance sets, since linear growth requires only
Gaussian sixth moments.

For a centered row product $X$, Taylor's remainder gives
$|\mathbb E(e^{itX/\sqrt n}\mid\mathcal F_n)
-1+t^2\mathbb E(X^2\mid\mathcal F_n)/(2n)|
\le |t|^3\mathbb E(|X|^3\mid\mathcal F_n)/(6n^{3/2})$.
Raise this identity to the $n$th power in each of the two independent
groups.  It proves that the conditional law of
$Z_n=\sqrt n(g_n'(0)-\mathbb E[g_n'(0)\mid\mathcal F_n])$
converges weakly in probability to $N(0,s_*^2)$.
To justify this random-law statement, every subsequence has a further
subsequence on which both covariance matrices converge almost surely;
the preceding characteristic-function calculation and the continuity
theorem then apply deterministically on almost every such outcome.
For a continuous limiting CDF, weak convergence is uniform in its
argument: choose finite grids with small Gaussian CDF increments and
small tails, and use monotonicity between grid points.  Therefore the
conditional Kolmogorov distance $d_n$ tends to zero in probability,
and $\mathbb E d_n\to0$ since $0\le d_n\le1$.

For any $\mathcal F_n$-measurable shift $b_n$, a centered Gaussian has
its largest interval probability at the interval centered at zero.
Consequently
\[
 \Pr\{|Z_n+b_n|\le us_*\mid\mathcal F_n\}
 \le2\Phi(u)-1+2d_n.
\]
The Gaussian assertion follows by differentiating the interval mass
with respect to its center and using the even, decreasing Gaussian
density.  Take
$b_n=\sqrt n\,\mathbb E[g_n'(0)\mid\mathcal F_n]$ and average.
This proves \eqref{cp:eq-scalar-fluctuation}; previous layers may
produce an arbitrary shift and cannot increase the limiting small-ball
probability.  No recursive matrix central limit theorem is needed.

\noindent\textbf{From the derivative to the trajectory.}\quad
Use the finite-query conclusion \eqref{cp:finite-time} of
\Cref{cp:source} for the two runs and the finite training query list.
Put $\lambda=\gamma/m$, $\ell=\log(en)$ and
$r=1/(\lambda\sqrt\ell)$. Each run is bounded there by
$16\beta^{6L}Y/\lambda$, so their difference has twice this bound.
The parameter-two ellipse of $[0,r]$ has real projection
$[-r/8,9r/8]$ and imaginary projection $[-3r/8,3r/8]$, and hence lies
inside the supplied rectangle.  On its event, $g_n$ is bounded by
$32\beta^{6L}Y/\lambda$.  Choose
$N=\lceil2\ell/\log2\rceil\le4\ell$.
Lemma~\ref{cp:derivative-transfer} gives
\[
 \sup_{0\le t\le r}|g_n(t)|
 \ge\frac{|g_n'(0)|}{32\lambda\ell^{5/2}}
       -\frac{768\beta^{6L}Y}{\lambda n^2}.
\]
For fixed $u>0$, the second term is eventually at most half the first
on $|g_n'(0)|\ge us_*/\sqrt n$.  The two source failures tend to zero;
no independence from the derivative event is required.  Thus
the conditional fluctuation claim gives a limiting lower probability
at least $2[1-\Phi(u)]$ for threshold
$us_*/(64\lambda\sqrt n\ell^{5/2})$.
Insert the lower bound on $s_*$, use $\lambda=\gamma/m$ and
$1/(64\sqrt2)\ge1/128$, and choose
$u=\Phi^{-1}(1/2+\delta/4)$.  The limiting probability is at least
$1-\delta/2$, leaving a strict margin to obtain $1-\delta$ eventually.
Since $g_n(0)=0$, a positive threshold on the compact witness interval
requires a strictly positive time.  That training query belongs to
each promised domain.  At its fitted endpoint the predictions equal
the same label, consistently with the asserted scope.
The asserted constant can therefore be chosen as the following local
quantity:
\[
 c_{\phi,L,\delta}=
 \frac{\Phi^{-1}(1/2+\delta/4)}{128}
 \sqrt{\frac{\mu_{2,L}}{\mu_4}}>0.
\]
Its moment ratio depends only on the activations and depth, as required.
\end{proof}
